\documentclass[11pt, a4paper]{article}
\title{Ingeniería Informatica (UIB)\\Estructura de Computadores I (2022-2023)\\Grupo 301}
\author{Serafí Nebot Ginard (41617163C)\\sng656@id.uib.eu}
\date{}
\thispagestyle{empty}

\usepackage[legalpaper, margin=2.5cm]{geometry}
\usepackage{hyperref}
\usepackage{fancyhdr}
\usepackage{amsmath}
\usepackage{amssymb}
\usepackage{parskip}
\usepackage{listings}
\usepackage{indentfirst}
\usepackage{titling}
\usepackage[flushleft]{threeparttable}
\usepackage{framed}
\usepackage[page,toc,titletoc,title]{appendix}
\renewcommand\maketitlehooka{\null\mbox{}\vfill}
\renewcommand\maketitlehookd{\vfill\null}

\lstdefinestyle{customasm}{
	belowcaptionskip=1\baselineskip,
	frame=single,
	xleftmargin=\parindent,
	language=[Motorola68k]Assembler,
	basicstyle=\footnotesize\ttfamily,
	captionpos=b
	% commentstyle=\color{green},
	% keywordstyle=\color{red}
}

\pagestyle{fancy}
\fancyfoot{}
\lhead{}
\rhead{}
\lfoot{Estructura de Computadores I}
\rfoot{\thepage}
\renewcommand{\headrulewidth}{0pt}
\renewcommand{\footrulewidth}{0.4pt}
\setlength{\headheight}{16pt}

\begin{document}

\renewcommand{\familydefault}{ppl}
\fontfamily{ppl}
\selectfont

\begin{titlingpage}
	\clearpage\maketitle
\end{titlingpage}
\thispagestyle{empty}
\pagebreak

\tableofcontents

\pagebreak
\section{Acronyms}
\textbf{MS[B,W]}: Most Singificant Byte/Word (Byte/Word whose address is greater)

\textbf{LS[B,W]}: Least Singificant Byte/Word (Byte/Word whose address is lower)

When referencing anything related to the HAL9000, the emulated machine, has been prepended with an
"e" to dinstinguish between the host machine and the emulated machine.

\pagebreak
\section{Reference tables}

\subsection{Subroutines}
Input and Output columns indicate the subroutine's stack layout before and after its execution. Each
item occupies 1 word (16 bits) and the first item is the value pointed by the \emph{Stack Pointer}
($\mathbf{SP}$).
\begin{table}[h!]
	\centering
	{
		\fontfamily{lmtt}
		\selectfont
		\begin{threeparttable}
			\begin{tabular}{ | l | l |  l | l | l | }
				\hline
				\textbf{Name} & \textbf{Description} & \textbf{Type} & \textbf{Input} & \textbf{Output} \\
				\hline
				RGET & Get value of eregister & library &
					\begin{tabular}[t]{ c }
						eregister number
					\end{tabular} &
					\begin{tabular}[t]{ c }
						eregister value
					\end{tabular} \\
				\hline
				RSET & Set value of eregister & library &
					\begin{tabular}[t]{ c }
						eregister value \\
						eregister number
					\end{tabular} &
					\begin{tabular}[t]{ c }
						eregister value \\
						eregister number
					\end{tabular} \\
				\hline
				FLAGS & Update specific eflags & library & 
					\begin{tabular}[t]{ c }
						value (MSW) \\
						value (LSW) \\
						selected eflags
					\end{tabular} &
					\begin{tabular}[t]{ c }
						value (MSW) \\
						value (LSW) \\
						selected eflags
					\end{tabular} \\
				\hline
			\end{tabular}
			\begin{tablenotes}
				\small
				\item 
			\end{tablenotes}
		\end{threeparttable}
	}
	\label{table:subroutines}
	\caption{Subroutines}
\end{table}

\subsection{Commonly used registers}
\begin{table}[h!]
	\centering
	{
		\fontfamily{lmtt}
		\selectfont
		\begin{tabular}{ | l | l | l | }
			\hline
			\textbf{Name} & \textbf{Phase} & \textbf{Description} \\
			\hline
			D2 & Execution & Instruction's first operand (if needed) \\
			\hline
			D3 & Execution & Instruction's second operand (if needed) \\
			\hline
			D4 & Execution & Instruction's third operand (if needed) \\
			\hline
			A2 & Execution & Instruction's first operand's ememory address (if needed) \\
			\hline
			A3 & Execution & Instruction's second operand's ememory address (if needed) \\
			\hline
			A4 & Execution & Instruction's third operand's ememory address (if needed) \\
			\hline
		\end{tabular}
	}
	\label{table:registers}
	\caption{Commonly used registers}
\end{table}

\pagebreak
\section{Fetch phase}
The fetch phase of the HAL9000 is fairly simple:
\begin{enumerate}
	\item Load 16 bits from ememory pointed by $\mathbf{EPC}$ into $\mathbf{EIR}$
	\item Increase $\mathbf{EPC}$ by 1
\end{enumerate}
Actually, it is not that simple. We have to take into account the fact that each memory cell of the HAL9000, unlike the Motorola
68000, holds 16 bits, instead of 8. This is why the $\mathbf{EPC}$ register is incremented by 1
instead of 2. Therefore, in order to load the instructions from ememory directly, the value of
$\mathbf{EPC}$ has to be multiplied by 2 ($16 / 8 = 2$).

So, the updated version of the previous steps is:
\begin{enumerate}
	\item Load 16 bits from ememory pointed by $\mathbf{EPC} \cdot 2$ into $\mathbf{EIR}$
	\item Increase $\mathbf{EPC}$ by 1
\end{enumerate}

This steps can be implemented like the following Motorola 68000 assembly:
\begin{lstlisting}[style=customasm, caption={HAL9000 emulator fetch phase code}, label={code:fetchcode}]
MOVE.W EPC, D1
MULU #2, D1		; multiply address by 2 (each position contains 2 Bytes)
MOVEA.W D1, A0
MOVE.W EMEM(A0), EIR	; load current instruction to EIR
ADDQ.W #1, EPC		; increase EPC by 1 after fetch
\end{lstlisting}

\pagebreak
\section{Decode phase}
\subsection{Decode phase introduction}

When developing an instruction decoder it is very common to use a Binary Decision Tree, a structure
based on a sequential decision process. This technique is very useful when the instructions'
opcodes to decode are completely arbitrary, with no common pattern between them and their assigned
ID. Luckily, this is not the case for the instruction set of the HAL9000 given for this assignment,
shown in Table \ref{table:halinstructions}, as each instruction's ID is encoded within its opcode
(with some exceptions).

\begin{table}[h]
	\centering
	{
		\fontfamily{lmtt}
		\selectfont
		\begin{tabular}{ | l | l | l | }
			\hline
			\textbf{ID} & \textbf{Instruction} & \textbf{Encoding} \\
			\hline
			0 & LOA M, Ti & 0000xxxxxxxxxxxx \\
			1 & STO Ti, M & 0001xxxxxxxxxxxx \\
			2 & LOIP (Xb), Ti & 0010xxxxxxxxxxxx \\
			3 & STIP Ti, (Xb) & 0011xxxxxxxxxxxx \\
			4 & GOI M & 0100xxxxxxxxxxxx \\
			5 & GOZ M & 0101xxxxxxxxxxxx \\
			6 & GON M & 0110xxxxxxxxxxxx \\
			7 & EXIT & 10xxxxxxxxxxxxxx \\
			8 & COPY Rb, R & 11000xxxxxxxxxxx \\
			9 & ADD Ra, Rb, R & 11001xxxxxxxxxxx \\
			10 & SUB Ra, Rb, Rc & 11010xxxxxxxxxxx \\
			11 & AND Ra, Rb, Rc & 11011xxxxxxxxxxx \\
			12 & SET \#k, Rc & 11100xxxxxxxxxxx \\
			13 & ADQ \#k, Rc & 11101xxxxxxxxxxx \\
			14 & LSH \#p, Rb, \#n & 11110xxxxxxxxxxx \\
			\hline
		\end{tabular}
	}
	\label{table:halinstructions}
	\caption{HAL9000 instruction set}
\end{table}

In order to take advantage of this pattern, let's separate the HAL9000 instruction set into two groups and
an exception.

\noindent\textbf{Group 0} \\
Instructions with ID from 0 to 6 follow the pattern:
\begin{equation*}
	\mathbf{0x_3x_2x_1k_{12}k_{11}k_{10}k_9k_8k_7k_6k_5k_4k_3k_2k_1}
\end{equation*}
where $\mathbf{x_{1 \ldots 3}}$ encode the instruction's ID as an unsigned binary number and
$\mathbf{k_{1 \ldots 12}}$ encode the rest of the instruction. This scenario is very easy to resolve
as the only thing needed to obtain the ID is to logically shift right the whole instruction by 12
bits:
\begin{equation*}
	\mathbf{0000000000000x_3x_2x_1}
\end{equation*}

\noindent\textbf{Group 1} \\
Instructions with ID from 8 to 14 follow the pattern:
\begin{equation*}
	\mathbf{1x_4x_3x_2x_1k_{11}k_{10}k_9k_8k_7k_6k_5k_4k_3k_2k_1}
\end{equation*}
where $\mathbf{x_{1 \ldots 4}}$ encode the instruction's ID as an unsigned binary number and
$\mathbf{k_{1 \ldots 11}}$ encode the rest of the instruction. This scenario may be resolved by
setting the MSb to 0 and logically shifting right the whole instruction by 11 bits:
\begin{equation*}
	\mathbf{000000000000x_4x_3x_2x_1}
\end{equation*}

\noindent\textbf{The exception} \\
Instruction with ID 7 is the only one which does not follow any pattern:
\begin{equation*}
	\mathbf{10k_{14}k_{13}k_{12}k_{11}k_{10}k_9k_8k_7k_6k_5k_4k_3k_2k_1}
\end{equation*}
where the two MSb ($\mathbf{10}$) is the instruction's opcode and $\mathbf{k_{1 \ldots 14}}$ encode
the rest of the instruction. This scenario may be resolved by checking the two MSb and hardcoding
the ID.

\pagebreak
\subsection{Decode phase implementation}
Due to the fact that the HAL9000 instructions can be divided into two groups, a one-level Binary Decision
Tree has been implemented in order to treat each group differently.

The decode algorithm has been implemented in the following manner:
\begin{enumerate}
	\item If the MSb is equal to $\mathbf{0}$ then the instruction is in \textbf{Group 0}. Go to
		step 6.
	\item If the MSb is equal to $\mathbf{1}$ then the instruction is in \textbf{Group 1} or
		\textbf{The exception}.
	\item Shift left by 1 bit to make the instruction's opcode the same size as \textbf{Group 0}'s.
	\item If the MSb is equal to $\mathbf{1}$ then the instruction is in \textbf{Group 1}. Got to
		step 6.
	\item Replace the instruction with $\mathbf{0111000000000000}$ (instruction is \textbf{The
		exception}; replace opcode with a hardcoded value).
	\item Shift right 12 bits to obtain the instruction's opcode.
\end{enumerate}
The decode algorithm can be translated into the following Motorola 68000 subroutine code:
\begin{lstlisting}[style=customasm, caption={HAL9000 instruction decoder subroutine}, label={code:h9kdecoder}]
DECOD:
	MOVE.L D0, -(SP)		; save current value of D0
	MOVE.W 8(SP), D0		; load instruction
	BTST.L #15, D0
	BEQ.S DECOD_END			; if bit 15 = 0 jump to end
	LSL.W #1, D0			; logically shift left to "normalize"
					; 5 bit opcode to 4 bit opcode
	BTST.L #15, D0
	BNE.S DECOD_END			; if bit 15 = 1 jump to end
	MOVE.W #%0111000000000000, D0	; if bit 15 = 0 then instruction = EXIT
DECOD_END:
	; logically shift right the four MSb to the four LSb
	LSR.W #8, D0
	LSR.W #4, D0
	MOVE.W D0, 10(SP)		; save instruction ID to reserved space
	MOVE.L (SP)+, D0		; restore value of D0
	RTS
\end{lstlisting}

Which can be used along with the following m68k code to decode the instruction in $\mathbf{ESR}$:
\begin{lstlisting}[style=customasm, caption={HAL9000 instruction decoder subroutine call}, label={code:h9kdecodercall}]
SUBQ.W #2, SP		; reserve 1 word for instruction id
MOVE.W EIR, -(SP)	; pass instruction to decode
JSR DECOD
ADDQ.W #2, SP		; pop instruction to decode from stack
MOVE.W (SP)+, D1	; retrieve instruction's id
\end{lstlisting}

\pagebreak
\section{Execution phase}
Each instruction has been implemented inside its own label so, each einstruction's execution code is
found at a different memory location (not all instruction execution occupy the same space!).

Therefore, a jump list has been used in order to jump to the appropiate einstruction's execution
code:
\begin{lstlisting}[style=customasm, caption={HAL9000 instruction execution jump list}, label={code:h9kexcjmplst}]
	; assuming D1 contains the einstruction's id
	MULU #6, D1 ; multiply by 6 as each JMP instruction occupies 6 Bytes
	MOVEA.L D1, A1
	JMP JMPLIST(A1)
JMPLIST:
	JMP ELOA
	JMP ESTO
	JMP ELOIP
	JMP ESTIP
	JMP EGOI
	JMP EGOZ
	JMP EGON
	JMP EEXIT
	JMP ECOPY
	JMP EADD
	JMP ESUB
	JMP EAND
	JMP ESET
	JMP EADQ
	JMP ELSH
\end{lstlisting}

$\mathbf{JMPLIST}$ is an ordered list of $\mathbf{JMP}$ instructions with the corresponding
einstruction's labels. For instance, if the current einstruction to execute is $\mathbf{EADD}$, the
value of $\mathbf{D1}$ will be 9, which is $\mathbf{EADD}$'s ID (see Table
\ref{table:halinstructions}). Then, $\mathbf{D1}$ will be multiplied by 6 to point to the 9'th entry
of the $\mathbf{JMPLIST}$.

\pagebreak
\appendix
\begin{appendices}
	\section{HAL9000 Compiler}
\end{appendices}

\end{document}

